\documentclass[11pt,a4paper]{article}
\usepackage[utf8]{inputenc}
\usepackage[T1]{fontenc}
\usepackage[spanish]{babel}
\usepackage{hyperref}
\usepackage{geometry}
\geometry{margin=2.5cm}
\usepackage{parskip}

\title{326 un absoluto que no se vuelve a mirar cobra siglos}
\author{Alejandro Valenzuela Rivera\\\small Universidad Autónoma de Nuevo León}
\date{\today}
\begin{document}
\maketitle

\begin{quote}
\textbf{Sesgo declarado.} Este trabajo se sostiene desde el sesgo declarado de su autor, y no finge neutralidad. Lo que afirma se puede tumbar, y ese es el punto: el falsificador está escrito, no escondido.
\end{quote}

\section*{Resumen}
Un crimen sigue teniendo consecuencias hoy no por su tamaño: porque se hizo sin pregunta. Lo que no se vuelve a mirar se vuelve institución.

\section{Por qué se sostiene}
La regla se sostiene mientras alguien la pueda contestar con lo que la esclavitud, y el abolicionismo muestra.

\section{De dónde viene}
Tesis/publicable/casos.md.

\section{El caso}
La esclavitud: se hizo sin pregunta, y la respuesta llegó dos siglos tarde.

\section{Método}
La pregunta que llega tarde encuentra el cambio ya encarnado en las constituciones, y no puede deshacerlo.

\section{Límites, y qué lo tumba}
Una condena rápida que no dejó rastro, y aun así se pudo preguntar antes.

\section{Lo que queda abierto}
\textquestiondown{}Hubo pregunta antes del acto, o doscientos años después?

\section*{Cómo se leyó esta tesis}
Esto se escribió en el repositorio man-made-empathy, del que este candado es una pieza: candados/06 del caso/326 un absoluto que no se vuelve a mirar cobra siglos.md. La máquina que lo compuso lo hizo desde este mismo archivo, y no agregó nada que no estuviera escrito. Si algo de esto es falso, se dice en el mismo sitio, y se corrige con fecha.

\end{document}
